\section{Calibrating the geometry to prices}\label{sec:calibration}

Theorem~\ref{thm:cal} asks for two guarantees from one scaling.
Elimination must produce a basis with bounded entries and bounded
expansion coefficients for unit vectors, and every
circuit must have a large coordinate wherever its assigned surplus is
maximal. The same matrix minors govern both requirements. We first
express them as comparisons of minor magnitudes, then show how to make
these comparisons using only short rational data.

\subsection{Cofactors for elimination and for circuits}

Let \(B_0\) be the given ordered basis, with \(r\) rows. Index its
nonzero minors by their row and column sets \(u\), and write \(\ell_u\)
for the base-two logarithm of the absolute determinant. For a list
\(T\) of \(j-1\) distinct columns, \(1\le j\le r\), define a row
vector \(c\) by
\begin{equation}\label{eq:cofactor}
 c_s=\det B_0\big[[j],(T,s)\big].
\end{equation}
The entry vanishes when \(s\) repeats a column of \(T\). Expansion
along the last column expresses \(c\) as a combination of the first
\(j\) rows of \(B_0\). The coefficient on row \(j\) is
\(\det B_0[[j-1],T]\), with the empty determinant equal to one.

There are two uses for this formula. In elimination, \(T\) is the
list of earlier pivots. Its determinant is nonzero, so \(c\) supplies
the next elimination row, with zeros at all earlier pivots. For a
circuit, take \(j=r\). The vectors of \(W\) supported on that circuit
form a line, so the columns outside its support have rank \(r-1\).
Choosing \(r-1\) independent columns there gives a cofactor vector
that vanishes outside the circuit. It is nonzero because \(B_0\) has
rank \(r\), and hence represents the circuit up to a scalar.

Consider for the moment the reciprocal scaling
\(D(p)_{i,h}=2^{-hp_i}\). Given a cofactor list \(T\) and a slot
\(s=(i,h)\) on its support, let \(u(s)\) be the associated minor
index. For a vector \(y\) of values on the minor indices, set
\begin{equation}\label{eq:rho}
 \rho_s(y,p)=y_{u(s)}-hp_i,\qquad
 \rho_s(\ell,p)=\log |(cD(p))_s|.
\end{equation}
Consequently a maximizing slot is a largest entry in absolute value,
and the difference from that maximum gives the logarithm of a
normalized entry magnitude.

The algorithm makes the comparisons using short rational substitutes
\(f_u\) for the logs. These substitutes need not have the same
numerical size as the logs. What we require is that for each query
\(p\) there exist prices \(\widehat p\) at which the true normalized
entry magnitudes agree, to the requested absolute accuracy, with
those computed from \((f,p)\). One vector \(\widehat p\) must work for
all the cofactor rows at that query, including circuits on which the
algorithm never requests a minor. Our construction will prove
existence of \(\widehat p\) by projecting a system of linear
inequalities.

\subsection{Linear types and online compression}

For two real vectors \(y,y'\) on the same finite index set, write
\(y\sim_M y'\) when every integer affine form of coefficient
\(1\)-norm at most \(M\), including its constant, has the same sign
in \(\{-1,0,1\}\) on \(y\) and \(y'\). This agreement of linear
types lets us transfer feasibility from one parameter vector to the
other. Compare the use of Frank--Tardos rounding in
\cite[Lemma~(6.6.1)]{GLS}.

\begin{lemma}[Projection]\label{lem:projection}
Consider a finite system of linear inequalities, strict or weak,
in \(k\) real unknowns and additional parameters \(y\).
Let all its affine forms have integer coefficients and
coefficient \(1\)-norm (including constant) at most \(A\), where \(A\ge2\).
Then feasibility is the same at any two parameter
vectors with \(y\sim_M y'\), where
\(M=(k+1) A^{k+1}\).

If there are no parameters and the system is feasible,
it has a solution with both a common positive denominator and the absolute
numerators at most \((A+1)^{k+1}\).
\end{lemma}
\begin{proof}
Fix the parameters. Give every strict inequality a common additional
margin \(\delta\) and write the resulting inequalities weakly. The
original system is feasible precisely when this weak system is
feasible for some \(\delta>0\). For fixed \(\delta\), Farkas' lemma
reduces feasibility to testing nonnegative combinations of the
inequalities that eliminate the \(k\) unknowns. It suffices to test
the extreme rays of this cone of multipliers.

On the support of an extreme ray the multiplier kernel has dimension
one. Its support thus has size at most \(k+1\), and subdeterminants
give an integral generator whose entries are at most \(A^k\) by
Hadamard's inequality. Each resulting affine form on the parameters,
with \(\delta\) set to zero, has coefficient norm at most
\((k+1)A^{k+1}\). Require the projected inequality to be strict if a
strict input inequality has a positive multiplier, and weak otherwise.
These conditions are necessary. They are also sufficient, since the
cone has finitely many extreme rays and one sufficiently small
\(\delta>0\) preserves every margin required by a strict combination.
Agreement of linear types preserves all these tests.

For the second assertion, let \(\delta\) be a new unknown and add
\(0\le\delta\le1\) to the system with margins. Intersect with a
closed orthant containing a solution with \(\delta>0\).
The new polyhedron is pointed and has a vertex maximizing \(\delta\).
Its constraint rows, with constants, have coefficient \(1\)-norm at
most \(A+1\). Cramer's rule and Hadamard's inequality now give the
claimed common denominator and numerator bounds.
\end{proof}

The next theorem, proved in Section~\ref{sec:compression}, furnishes
substitutes that agree in linear type with a small perturbation of
the true logs. A price query will use only the entries requested
during elimination. The guarantee for all the remaining minors is
what makes the conclusion about unseen circuits possible.

\begin{theorem}[Online compression]\label{thm:compression}
Given an \(r\times d\) real matrix (\(0\le r\le d\), \(d\ge1\))
and \(K\ge1\) an integer,
index its nonzero minors by \(u\) (specifying row and column sets)
and let
\(\ell_u\) be the base-two logarithm of their absolute values.
There is a deterministic procedure that, on request of an
individual index \(u\), supplies a rational \(f_u\)
or declares a restart.

At any time the answers since the most recent restart (if no restart,
all answers) admit a common assignment \(f\) of rationals to
\emph{all} the indices with a common positive denominator and absolute
numerators each bounded by \(F=2^{\poly(d,K)}\), such that
\begin{equation}\label{eq:compress}
 f\sim_{2^K} g\quad\text{for some }g\text{ with }
 \|g-\ell\|_\infty\le2^{-K}.
\end{equation}
The total number of restarts is bounded by a polynomial in \(d\).
Initialization and each request (including responses with explicit
bits or restarts) take strongly polynomial time,
counting the value of \(K\) as additional input size. The integer bound \(F\)
is available at initialization.
\end{theorem}

\subsection{The common lift}

For \(P\) in Theorem~\ref{thm:cal}, set
\[
 E=P+6,\qquad e=2^{-E},\qquad A=(E+5)2^E
\]
and take \(K\) least with \(2^K\ge(N+1)A^{N+1}\).
Fix price data \(p\) and assignments \(f,g\) for \(B_0\) as in
\eqref{eq:compress}. We encode the normalized cofactor magnitudes at
\((f,p)\) as follows. For each nonzero cofactor vector take the slot
\(t\) maximizing \(\rho_t(f,p)\), using a fixed tie rule. For each
slot \(s\) on the same support write
\[
 w=\rho_s(f,p)-\rho_t(f,p)\le0,\qquad
 \Delta=\rho_s(y,\widehat p)-\rho_t(y,\widehat p).
\]
In \(\Delta\) the vector \(y\) is the parameter and the \(N\) components
of \(\widehat p\) are the unknowns. Impose
\[
 \begin{cases}
 e\lfloor w/e\rfloor\le\Delta\le e\lceil w/e\rceil,
      & w\ge-E,\\
 \Delta\le-E, & w<-E.
 \end{cases}
\]
For a ratio smaller than \(2^{-E}\) it suffices to preserve that upper
bound. Ratios above the cutoff need only one grid cell to specify
their log. All constraints hold at \((y,\widehat p)=(f,p)\).
Multiplying by \(2^E\) gives integer forms of coefficient norm at
most \(A\), independent of the bit lengths of the original prices
and the proxies. Lemma~\ref{lem:projection} therefore gives a solution
\(\widehat p\) at \(y=g\), since \(f\sim_{2^K}g\). The system
includes every cofactor list, so the solution is a single price vector
shared by all of them. The algorithm never forms this system.

Replacing \(g\) by the true log vector \(\ell\) changes each
difference by at most \(2\cdot2^{-K}\le2e\). Thus at the lifted
prices the true difference is within \(3e\) of \(w\) for
\(w\ge-E\), and at most \(-E+2e\) for \(w<-E\). Taking
\(D=D(\widehat p)\) yields
\begin{equation}\label{eq:pivot-estimate}
 \left|\frac{|(cD)_s|}{|(cD)_t|}-2^w\right|\le8e.
\end{equation}
In the first case this follows from the mean value theorem applied to
\(2^x\) for \(x\le3e\). In the second, both magnitudes are at most
\(2^{-E+2e}\).

\subsection{The calibration procedure}

Use Theorem~\ref{thm:compression}
on \(B_0\), with \(K\) as just specified. Take \(R=4F\).
To answer a price query \(p\), grow a sequence of pivot slots \(T\), starting empty.
At row \(j=1,\ldots,r\), take the cofactor vector with
\(T\) consisting of the previous pivots.
To see there will be a pivot available, recall that this vector \(c\)
is in the span of the first \(j\) rows. Its coefficient on original row \(j\) is \(d_{j-1}=\det B_0[[j-1],T]\), nonzero
by construction (\(d_0=1\)). Thus \(c\ne0\).
Request \(f\) values for all its nonzero entries.
If a restart occurs, abandon the query with that report.
Otherwise compute the row as follows.
\begin{itemize}
 \item As pivot at this row take slot \(t\) maximizing
 \(\rho_s(f,p)\) on the support, breaking ties by the rule fixed above.
 This gives \(d_j=c_t\ne0\) for the next step.
 \item For row \(j\) of the answer \(B\),
 keep zeros, and at each supported slot \(s\) put a rational
 within \(e\) of
 \[
    \sgn(d_{j-1}c_s)\cdot 2^{\rho_s(f,p)-\rho_t(f,p)}.
 \]
\end{itemize}
One query visits at most \(r\) cofactor rows and requests at most
\(2Nr\) minor proxies. Signs and zeros use minors of the original
matrix, computed as in Section~\ref{sec:arithmetic}.
Theorem~\ref{thm:compression} bounds the cost of the requests and
the total number of restarts.
The rational approximations take polynomial bit time.
For exponents below \(-E\), use zero.
Otherwise, split off the integer power, leaving \(2^q\) with
\(0\le q<1\).
Use \(\ln 2=\sum_{j\ge1}2^{-j}/j\) and then the series for
\(\exp(q\ln 2)\), e.g., through \(E+10\) in each,
by geometric tail bounds.

\begin{proof}[Completion of the proof of Theorem~\ref{thm:cal} given Theorem~\ref{thm:compression}]
Fix any uninterrupted sequence of answers from initialization or the
last restart up to the present time. Theorem~\ref{thm:compression}
extends all the returned proxies to one vector \(f\). For each query
in the sequence choose the scaling \(D\) from its common lift.
For row \(j\) of \(B_*\), use the corresponding \(cD\)
with sign adjusted by \(\sgn(d_{j-1})\)
and scaled to absolute entry one at the pivot.
Thus it now has a positive coefficient on row \(j\)
of \(B_0D\), giving \(B_*=L B_0D\)
with \(L\) as required.
By~\eqref{eq:pivot-estimate}, \(B\) is within \(9e\le2^{-P}\)
entrywise, and \(B_*\)
has entries bounded by two in absolute value.
Any vector of sup norm one in \(WD\)
can be expanded with coefficient on row \(j\)
bounded by \(3^{j-1}\le2^{100N^2}\)
in absolute value. Indeed, subtract row multiples in pivot order.
The pivot entry has magnitude one, so a residual of norm \(M\) requires
a coefficient of magnitude at most \(M\). The row itself has norm at
most two, leaving a residual of norm at most \(3M\).

For each circuit choose one cofactor representative and assign demands
\[
 a_s = 1+F+\lfloor f_{u(s)}\rfloor
\]
on its support. Since \(|f_{u(s)}|\le F\), these demands belong to
\([1,1+2F]\subseteq[1,R]\). They use the same \(f\) at every query in
the sequence.
Fix a circuit, with \(t\) as in the common lift
(so maximizing \(\rho_s(f,p)\) for this representative).
If \(s=(i,h)\) maximizes surplus \(a_s-hp_i\), then
\(w=\rho_s(f,p)-\rho_t(f,p)>-1\),
by the bound of one on rounding loss.
So in the common lift,
\(\rho_s(\ell,\widehat p)-\rho_t(\ell,\widehat p)\)
is greater than \(-1-3e\).
At any slot on this support, the true scaled log difference
from \(t\) is at most \(3e\).
Thus after scaling the circuit by \(D\), the entry at \(s\)
has magnitude at least
\(2^{-1-6e}\ge1/3\) times the sup norm.
\end{proof}
